% !TeX root = ../../main.tex
\subsection{助動詞 shall}
    \settowidth{\dimen0}{Shall I ～？}
    \makebox[\dimen0][l]{Shall I ～？}「～ \fitblankbf[]{しましょうか} ？」
    \vspace{0.5em}

    Shall we ～？ = \fitblankbf[]{Let's} + 動詞の原形 = \fitblankbf[]{Why don't we} 〜 ?
    \vspace{0.5em}

    \makebox[\dimen0][l]{}「～ \fitblankbf[]{しませんか} ？」

    % \begin{enumerate}
    %     \item Shall I open the window?
    %     \dialogueblank{窓を開けましょうか。}
    %     \vspace{0.5em}
    %     \item Shall we go to the park?
    %     \dialogueblank{公園に行きませんか。}
    %     \vspace{0.5em}
    %     \item Why don't we study together?
    %     \dialogueblank{一緒に勉強しませんか。}
    % \end{enumerate}

\subsection{would like to}
    \settowidth{\dimen0}{would like to + 動詞の原形 = \fitblankbf[]{want to} + 動詞の原形 }
    \makebox[\dimen0][l]{would like to + 動詞の原形 = \fitblankbf[]{want to} + 動詞の原形 }
    「〜 \fitblankbf[]{したい} 」
    \vspace{0.5em}

    ※ want 単体で使う時の意味としては 「～を \fitblankbf[]{ほしいと思う} 」

    % \begin{enumerate}
    %     \item I would like to visit Hokkaido.
    %     \dialogueblank{私は北海道を訪れたい。}
    %     \vspace{0.5em}
    %     \item She wants a new bag.
    %     \dialogueblank{彼女は新しいかばんがほしいと思う。}
    % \end{enumerate}

\subsection{would you like}
    \settowidth{\dimen0}{--- No , thank you . }
    \addtolength{\dimen0}{2em}
    \makebox[\dimen0][l]{would you like ～？}「～は \fitblankbf[]{いかがですか} ？」
    \vspace{0.5em}

    \makebox[\dimen0][l]{\hspace{2em}--- Yes , please . }「はい、 \fitblankbf[]{お願いします。} 」
    \vspace{0.5em}

    \hspace{2em}--- No , thank you . 「いいえ、 \fitblankbf[]{結構です。} 」

    % \begin{enumerate}
    %     \item Would you like some tea?
    %     \dialogueblank{お茶はいかがですか。}
    %     \vspace{0.5em}
    %     \item --- Would you like to join us?\\
    %     \hspace{2em}--- Yes, please.
    %     \dialogueblank{--- 私たちに加わりませんか。\\
    %     \hspace{2em}--- はい、お願いします。}
    %     \vspace{0.5em}
    %     \item --- Would you like some cake?\\
    %     \hspace{2em}--- No, thank you.
    %     \dialogueblank{--- ケーキはいかがですか。\\
    %     \hspace{2em}--- いいえ、結構です。}
    % \end{enumerate}
